\chapter{为什么程序字节还不是可运行映像}

\begin{statebox}
\textbf{项目里程碑 v0.4。} Stage 1 已进入长模式并拥有覆盖装载窗口的临时页表。
现在要从磁盘读取 Kernel 容器，先验证所有 PT\_LOAD 的范围、权限和相互关系，再复制段、
清零 BSS 并交接 BootInfo。项目解决的是第一次建立内核；深入机制再把同一事务扩展为
已有 task 的 \code{exec} 映像替换。已完成的 v0.8 复用这条规则校验用户 ELF：
限制可装载段和总页数、拒绝 W+X 或重叠区间，并要求入口属于可执行段。最新整机
验证已经采用三份用户 ELF，并证明截断映像会在进入 Ring 3 前被拒绝；这完成了
最小用户映像事务，但尚未扩展为具有参数、文件描述符和多进程地址空间的
\code{exec}。
\end{statebox}

\inputchapterbody{chapters/ch10-exec-elf-loader}
\inputdetail{chapters18/expansions/ch04-exec-transaction}

\begin{problemframe}
\textbf{尚未解决的问题。} ELF 映像和 BSS 已经建立，BootInfo 也到达入口，但内核不能
长期依赖 Stage 1 的 GDT、栈和异常状态。下一章从项目 v0.5 的描述符接管与统一异常帧
开始，解释 v0.8 用户请求实际采用的受控系统调用入口。
\end{problemframe}
